%------------------------------------------------------------------------------%
\section{Cálculo da Série de Taylor para o Arco Tangente}

Vamos ilustrar uma forma alternativa para calcular a Série de Taylor da
função $\arctan(x)$ para ilustrar algumas manipulações possíveis
quando precisamos obter uma Série de Taylor.
Vamos começar, como em outros casos, calculando as derivadas da função
\begin{align*}
  \frac{d}{dx}     & \arctan(x) = \frac{1}{1+x^2}                     \\[3mm]
  \frac{d^2}{dx^2} & \arctan(x) = \frac{-2x}{(1+x^2)^2}               \\[3mm]
  \frac{d^3}{dx^3} & \arctan(x) = \frac{6x^2-2}{(1+x^2)^3}            \\[3mm]
  \frac{d^4}{dx^4} & \arctan(x) = \frac{24x(1-x^2)}{(1+x^2)^4}   \;.
\end{align*}
Não precisamos ir mais longe para perceber que esse não é um bom método,
as expressões estão se tornando cada vez mais complexas,
portanto, precisamos tentar outra abordagem.

Inicialmente, nesse desenvolvimento, não vamos nos preocupar com o
rigor matemático para obter uma candidata a Série de Taylor,
depois justificamos o processo utilizado demonstrando a
convergência da série.
Começamos observando que a derivada da função arco tangente
é igual a soma de uma Série Geométrica com $\alpha=1$ e $r=-x^2$
\[
\frac{d}{dx} \arctan(x)
= \frac{1}{1+x^2}
= \frac{\alpha}{1-r}
= 1 - x^2 + x^4 - x^6 + x^8 - \cdots
\]
Escrevemos então que
\[
\frac{d}{dx} \arctan(x) =  1 - x^2 + x^4 - x^6 + x^8 - \cdots
\]
integrando os dois lados temos a candidata a Série de Taylor para $\arctan(x)$
\[
\arctan(x) = x -\frac{x^3}{3} + \frac{x^5}{5} - \frac{x^7}{7} + \frac{x^9}{9} - \cdots
\]
Temos uma Série de Taylor, porém, nossos cálculos não foram cuidadosos.
Precisamos demonstrar que a série obtida converge para a função.
Faremos isso voltando para a fórmula da soma da Série Geométrica
\begin{equation}
  \label{eq:arctan_serie_geometrica}
  \frac{1}{1+x^2}
  = 1 - x^2 + x^4 - x^6 + \cdots
  + (-1)^n x^{2n}
  + (-1)^{n+1} x^{2(n+1)} + \cdots
\end{equation}
Note que para qualquer $n$ o resto dessa série, $R_n$,
é uma Série Geométrica com $r=-x^2$ e $\alpha=(-1)^{n+1} x^{2(n+1)}$
\[
R_n =
(-1)^{n+1} x^{2(n+1)} +
(-1)^{n+2} x^{2(n+2)} +
(-1)^{n+3} x^{2(n+3)} + \cdots
\]
Utilizando a fórmula da soma dos termos dessa série obtemos
\[
R_n = \frac{(-1)^{n+1} x^{2(n+1)}}{1+x^2}    \;.
\]
Podemos então escrever a série em \eqref{eq:arctan_serie_geometrica}
como uma soma finita
\[
\frac{1}{1+x^2}
= 1 - x^2 + x^4 - x^6 + \cdots
+ (-1)^n x^{2n}
+ \frac{(-1)^{n+1} x^{2(n+1)}}{1+x^2}
\]
Temos agora uma expressão finita para a derivada da função arco tangente
\[
\frac{d}{dx} \arctan(x)
= 1 - x^2 + x^4 - x^6
+ \cdots
+ (-1)^n x^{2n}
+ \frac{(-1)^{n+1} x^{2(n+1)}}{1+x^2}
\]
Por ser uma soma finita podemos integrar termo a termo
sem preocupações com a convergência
\[
\arctan(x) =
x - \frac{x^3}{3}
+ \frac{x^5}{5}
- \frac{x^7}{7}
+ \cdots
+ \frac{(-1)^n x^{2n+1}}{2n+1}
+ \int_0^x \frac{(-1)^{n+1} u^{2(n+1)}}{1+u^2} \,du   \;.
\]
Estamos agora em uma situação similar ao Teorema de Taylor~\ref{teo:teorema-taylor},
onde podemos concluir que a série converge para a função se o erro
\[
R_n(x) = \int_0^x \frac{(-1)^{n+1} u^{2(n+1)}}{1+u^2} \,du
\]
convergir para zero, quando $n\to\infty$.

Para garantirmos que $R_n$ tende para zero, basta verificar que seu módulo tente a zero,
ou seja, queremos calcular o limite de
\begin{align*}
  \abs{R_n(x)}
  & = \abs{\int_0^x \frac{(-1)^{n+1} u^{2(n+1)}}{1+u^2} \,du} \\[3mm]
  & = \int_0^{\abs{x}} \abs{\frac{(-1)^{n+1} u^{2(n+1)}}{1+u^2}} \,du \\[3mm]
  & = \int_0^{\abs{x}} \frac{u^{2(n+1)}}{1+u^2} \,du
\end{align*}
quando $n\to\infty$. Como $1+u^2\geq 1$ temos que
\[
\abs{R_n(x)}
\leq \int_0^{\abs{x}} u^{2(n+1)} \,du
= \frac{u^{2n+3}}{2n+3} \evalat{0}{\abs{x}}
= \frac{\abs{x}^{2n+3}}{2n+3}    \;.
\]
Portanto, para todo $\abs{x} \leq 1$ sabemos que
\[
\lim_{n\to\infty} R_n(x) = 0
\]
e podemos concluir que, para todo $x\in[-1,\,1\,]$ a Série de Taylor da função
arco tangente converge para a função
\[
\arctan(x) = x - \frac{x^3}{3}
+ \frac{x^5}{5}
- \frac{x^7}{7}
+ \frac{x^9}{9}
- \cdots
\]

